\documentclass[10pt,a4paper]{amsart}
\usepackage[margin=19mm]{geometry}
\usepackage{amsmath,amssymb,mathtools}
\usepackage[scr=rsfs,cal=euler]{mathalfa}
\usepackage{xfrac}
\usepackage[unicode,hidelinks,bookmarks=false]{hyperref}
\newcommand{\ie}{i.e.\ }
\newcommand{\cf}{cf.\ }
\newcommand{\resp}{respectively\ }
\newcommand{\Subsection}{\subsection}
\providecommand{\nobreakdash}{\nobreak\hbox{-}}
\setlength{\emergencystretch}{2em}
\multlinegap0pt
\usepackage{tcolorbox}
\usepackage{amssymb}
\usepackage{mathtools}
\usepackage{dsfont}
\usepackage{float}
\usepackage{bm}
\usepackage{tikz}
\newtheorem{theorem}{Theorem}
\def\B{\mathscr{B}}
\def\Ce{\mathds{C}}
\def\D{\mathscr{D}}
\def\M{\mathscr{M}}
\def\Re{\mathds{R}}
\def\mod#1{|\mskip 1mu #1 \mskip 1mu|}
\def\modc#1{
    \left|
    \mskip 2mu #1 \vphan \mskip 2mu
    \right|
    }
\def\norm#1{\| \mskip 1mu #1 \mskip 1mu \|}
\def\normc#1{
    \left\|
    \mskip 2mu #1 \vphan \mskip 2mu
    \right\|
    }
\def\ran#1{\mathop{\mathrm{ran}}#1}
\def\ud{\,\mathrm{d}}
\def\vphan{\vphantom{\bigl|}}
\title{Interpolating between Tikhonov regularization and spectral cutoff}
\author{Martin Sæbye Carøe, Mirza Karamehmedovi\'c, Pierre Maréchal}
\date{Source: arXiv:2602.09505v1 (CC BY 4.0)}
\begin{document}
\maketitle

\noindent\emph{Source and licence.} Interpolating between Tikhonov regularization and spectral cutoff. Version arXiv:2602.09505v1 (CC BY 4.0), distributed by arXiv under the Creative Commons Attribution 4.0 International licence, http://creativecommons.org/licenses/by/4.0/, as recorded in the arXiv licence metadata for this exact version and shown on the abstract page https://arxiv.org/abs/2602.09505v1. Full licence notice: \texttt{licenses/arxiv-2602.09505-CC-BY-4.0.txt}.\par\smallskip

\noindent\textit{Editorial scope.} Counted unit: the article's theorem spectral-reg (source0.tex lines 409--433). The article's complete original proof of it is delivered with it (source0.tex lines 435--490, one \texttt{proof} environment of 56 lines). No mathematics is written, repaired or re-derived: the statement and the proof are the article's own lines, in the article's own notation, and no retained line is substituted or re-flowed.  Retained uncounted as the source-local context the proof needs: lines 335--351.  Nothing was omitted from any retained span, so the package carries no omission record.  No retained line contains a citation, so the wrapper carries no bibliography and no bibliography entry.  No retained line contains a figure environment, an image-inclusion command or drawing code, so no image is delivered and the package declares no asset dependency.  Editorial formatting only, in addition: an AMS \texttt{amsart} preamble that restates the article's own declarations and macros used by the retained lines, namely the environments \texttt{theorem} on separate counters, together with the standard packages those lines need.  The retained environments print in this document as Theorem 1 in order of appearance, whereas the article prints the same items with the numbering of its own full document.  The complete native source of the article as distributed in the arXiv e-print for this exact version is bundled unchanged as \texttt{sources/194321/source0.tex}.
\begin{theorem}\sf
\label{sve}
Let $T\colon F\to G$ be an injective bounded linear operator, where $F$ and $G$
are real Hilbert spaces. There then exist
\begin{enumerate}
\item
a Borel space $(\M,\B,\mu)$ with completely separable topology,
\item
an unitary operator $V\colon L^2(\M,\B,\mu)\to F$ and
an isometry $W\colon L^2(\M,\B,\mu)\to G$,
\item
an essentially bounded measurable function $\sigma\colon\M\to\Re$
that is strictly positive $\mu$-almost everywhere,
\end{enumerate}
such that $T=W[\sigma]V^*$, in which $[\sigma]$ denotes the operator of
multiplication by~$\sigma$. Moreover, $\ran{W}=\overline{\ran{T}}$. 
\end{theorem}

\begin{theorem}\sf
\label{spectral-reg}
Let~$T\colon F\to G$ be an injective bounded operator with
singular value expansion $T=W[\sigma]V^*$ as in Theorem~\ref{sve}.
Assume that $\sigma$ is not bounded away from zero, so that the
inverse problem $Tf=g$ is ill-posed.
Let $q\colon\Re_+^*\times\Re_+\to\Ce$ be such that
\begin{enumerate}
\item[(A1)]
for every $(\alpha,\sigma)\in\Re_+^*\times\Re_+$, 
$\mod{q(\alpha,\sigma)}\leq M$ for some positive constant~$M$,
\item[(A2)]
for every $\sigma\in\Re_+$, $\mod{q(\alpha,\sigma)}\leq c_\alpha\sigma$,
in which $c_\alpha$ is some positive constant.
\end{enumerate}
Then $R_\alpha:= V[\sigma^{-1}q(\alpha,\sigma)]W^*$ is well-defined and bounded,
and $\norm{R_\alpha}\leq c_\alpha$.
Moreover, if
\begin{enumerate}
\item[(A3)]
for every $\sigma\in\Re_+$,
$\lim_{\alpha\downarrow 0}q(\alpha,\sigma)=1$,
\end{enumerate}
then $(R_\alpha)$ is a regularization of~$T^\dagger$.
\end{theorem}

\begin{proof}
Let~$g\in\D(T^\dagger)$ and let $\bar{g}$ be its projection onto the
closure of the range of~$T$. Observe that, since~$g$
belongs to the domain of~$T^\dagger$, its projection~$\bar{g}$
must actually belong to the range of~$T$. Hence $\bar{g}=Tf$
for some $f\in F$, from which we deduce that
$$
[\sigma^{-1}]W^*\bar{g}=
[\sigma^{-1}]W^*W[\sigma]V^*f=
V^*f\in L^2(\M,\B,\mu).
$$
By Assumption (A1), $[q(\alpha,\sigma)][\sigma^{-1}]W^*g=[q(\alpha,\sigma)]V^*f$
also belongs to $L^2(\M,\B,\mu)$, which in turn implies that
$$
V[q(\alpha,\sigma)][\sigma^{-1}]W^*g=:R_\alpha g
$$
is well-defined.
Under Assumption (A2), we have, for every $g\in G$,
\begin{eqnarray*}
\normc{R_\alpha g}^2
&=&
\normc{V[\sigma^{-1}q(\alpha,\sigma)]W^*g}^2\\
&=&
\normc{[\sigma^{-1}q(\alpha,\sigma)]W^*g}^2\\
&\leq&
\normc{\sigma^{-1}q(\alpha,\sigma)}_\infty^2\normc{W^*g}^2\\
&=&
\normc{\sigma^{-1}q(\alpha,\sigma)}_\infty^2\normc{g}^2\\
&\leq&
c_\alpha^2\normc{g}^2,
\end{eqnarray*}
in which the second equality is due to the unitarity of~$V$
and the third equality is due to the fact that~$W$ is an isometry
(so that $W^*W=I$, where~$I$ denotes the identity).
The first assertion in the theorem follows.

Next, let $f^\dagger:= T^\dagger g$ and $f_\alpha:=R_\alpha g$.
From the definition of~$R_\alpha$ and the formula for~$T^\dagger$,
we readily see that
$$
f^\dagger-f_\alpha=
V\big[\sigma^{-1}(1-q(\alpha,\sigma))\big]W^*g=
V\big[1-q(\alpha,\sigma)\big][\sigma^{-1}]W^*\bar{g}=
V\big[1-q(\alpha,\sigma)\big]V^*f.
$$
Therefore, using the unitarity of~$V$, 
$$
\normc{f^\dagger-f_\alpha}^2=
\normc{\big[1-q(\alpha,\sigma)\big]V^*f}_{L^2(\M,\B,\mu)}^2=
\int\modc{1-q(\alpha,\sigma)}^2 \big(V^*f(\sigma)\big)^2\ud\mu(\sigma).
$$
Finally, under Assumption (A3), the function
$\sigma\mapsto \modc{1-q(\alpha,\sigma)}^2$ converges pointwise
to zero, and Lebesgue's dominated convergence theorem then implies
that $f_\alpha\to f^\dagger$ as $\alpha\downarrow 0$.
\end{proof}

\end{document}
